\fancyhead[C]{\textit{Solutions}}

\section{Solutions to Section 1}
	\begin{table}[H]
		\centering
		\begin{tabular}{c@{\hspace{2.5em}}cp{0.75\textwidth}}
			true & false & \\
			\specialrule{\lightrulewidth}{0.25em}{1.25em}
			\dbubble & \bubble & If \(A \in \mathbb{R}^{n \times n}\) has \(n\) linearly independent eigenvectors, then so does \(A^T\). \\[0.75em]
			\bubble & \dbubble & It is possible that there exists a square matrix \(A\) that is similar to a diagonalizable matrix \(B\), but \(A\) isn't diagonalizable. \\[0.75em]
			\bubble & \dbubble & Row-equivalent matrices are similar. \\[0.75em]
			\bubble & \dbubble & Every invertible matrix is diagonalizable.\\[0.75em]
			\bottomrule
		\end{tabular}
	\end{table}

\subsection{Justifications}
\paragraph{Justification to \textit{Q1}} If \(A\) has \(n\) linearly independent eigenvectors, then \(A\) could be diagonalized.
		\[A = PDP^{-1} \implies A^{T} = (PDP^{-1})^{T} = (P^{-1})^{T}DP^{T}.\] Since \(D\) is diagonal, the columns of \((P^{-1})^{T}\) 
		\emph{necessarily} form a set of eigenvectors of \(A^{T}\). (To know why, follow the proof provided in the book.) We know they are linearly independent
		because the columns form an invertible matrix.

\paragraph{Justification to \textit{Q2}} If \(A\) is similar to a matrix \(B\), then \(A = PBP^{-1}\) for some invertible matrix \(P\). If \(B\) is diagonalizable, then 
		\(B = QDQ^{-1}\), for some invertible matrix \(Q\), and a diagonal matrix \(D\). Therefore, 
		\[A = PBP^{-1} \implies A = P(QDQ^{-1})P^{-1} \implies A = (PQ)D(PQ)^{-1}.\] It therefore follows that, by definition, \(A\) \emph{should} be 
		invertible. Hence, the statement is false.

\paragraph{Justification to \textit{Q3}} Consider \(A = \begin{pmatrix}1 & 0 \\ 0 & 5\end{pmatrix}\) and \(A = \begin{pmatrix}1 & 0 \\ 0 & 1\end{pmatrix}\). Those
		two matrices are row equivalent but are certainly not similar.

\paragraph{Justification to \textit{Q4}} Consider \(A = \begin{pmatrix}1 & 1 \\ 0 & 1\end{pmatrix}\).


\section{Solutions to Section 2}
	\begin{table}[H]
		\centering
		\begin{tabular}{c@{\hspace{2.5em}}cp{0.75\textwidth}}
			possible & impossible & \\
			\specialrule{\lightrulewidth}{0.25em}{1.25em}
			\bubble & \dbubble & A real \(5 \times 5\) matrix that has 3 complex eigenvalues and 2 real eigenvalues. \\[0.75em]
			\bubble & \dbubble & A real \(2 \times 2\) matrix that has \(1 + 5i\) and \(-1 + 5i\) as eigenvalues. \\[0.75em]
			\dbubble & \bubble & A real \(3 \times 3\) matrix that has a real eigenvalue and 2 complex eigenvalues. \\[0.75em]
			\bottomrule
		\end{tabular}
	\end{table}
	
	\subsection{Justifications}
	\paragraph{Justification to Q1} For a real, square matrix, we have it that the complex eigenvalues come in conjugate pairs. Meaning, it is possible to have an odd number of
			complex eigenvalues if the matrix where to be real.

	\paragraph{Justification to Q2} \(1 + 5i\) and \(1 - 5i\) are \emph{not} complex conjugates.

	\paragraph{Justification to Q3} Consider \(A = \begin{pmatrix}1 & 2 & 0 \\ -2 & 1 & 0 \\ 0 & 0 & 1\end{pmatrix}\).

	\subsection{Solution to Second Question}
	We automatically know that the second eigenvalue is \(2 - i\). We now have to find a basis for each eigenspace. We use the usual method. For \(\lambda_1 = 2 + i\):
	\[\begin{pmatrix}2 - (2 + i) & -1 \\ 1 & 2 - (2 + i)\end{pmatrix} = \begin{pmatrix}-i & -1 \\ 1 & -i\end{pmatrix} \sim \begin{pmatrix} 1 & -i \\ 0 & 0\end{pmatrix}.\] 

	So the paramteric vector form of the solution set is \(\vec{x} = t\begin{pmatrix}i \\ 1\end{pmatrix}\). Choose \(\vec{v} = \begin{pmatrix}i \\ 1\end{pmatrix}\).

	We know that the conjugate of a complex eigenvector is also an eigenvector of a real matrix, so the corresponding paramteric vector form for the solution set is 
	\(\vec{y} = t\begin{pmatrix}-i \\ 1\end{pmatrix}\). Choose \(\vec{v}_2 = \begin{pmatrix}-i \\ 1\end{pmatrix}\).